%================================================
% ONBOARDING GUIDE — SUVA PROTOCOL (ENGLISH)
% Post-Quantum Signatures on EVM — Beginner Guide
% Target: 2nd-year student, basic programming knowledge
% Author: Mohamed Lamine MBENGUE
% Date: April 2026
%================================================
\documentclass[11pt,a4paper]{article}

\usepackage[utf8]{inputenc}
\usepackage[T1]{fontenc}
\usepackage[english]{babel}
\usepackage{geometry}
\geometry{margin=2.5cm}

\usepackage{amsmath,amssymb,amsthm}
\usepackage{booktabs}
\usepackage[table]{xcolor}
\usepackage{listings}
\usepackage{hyperref}
\usepackage{graphicx}
\usepackage{tikz}
\usetikzlibrary{arrows.meta, positioning, shapes.geometric, calc}
\usepackage{fancyhdr}
\usepackage{float}
\usepackage{array}
\usepackage{multirow}
\usepackage{longtable}
\usepackage{tcolorbox}
\tcbuselibrary{skins, breakable}

%--- Colours ---
\definecolor{suva-blue}{RGB}{0, 82, 165}
\definecolor{suva-green}{RGB}{0, 150, 90}
\definecolor{suva-orange}{RGB}{220, 100, 0}
\definecolor{code-bg}{RGB}{245, 245, 245}
\definecolor{pass-green}{RGB}{0, 140, 70}
\definecolor{fail-red}{RGB}{180, 0, 0}

%--- Header / Footer ---
\pagestyle{fancy}
\fancyhf{}
\rhead{\textcolor{suva-blue}{\textbf{SuVa Protocol --- Onboarding Guide}}}
\lhead{\textit{English version}}
\rfoot{\thepage}
\lfoot{\textit{April 2026}}

%--- Code style: Solidity ---
\lstdefinestyle{solidity}{
  language=C,
  basicstyle=\ttfamily\small,
  backgroundcolor=\color{code-bg},
  keywordstyle=\color{suva-blue}\bfseries,
  commentstyle=\color{gray}\itshape,
  numbers=left,
  numberstyle=\tiny\color{gray},
  frame=single,
  breaklines=true,
  captionpos=b
}

%--- Code style: Python ---
\lstdefinestyle{python}{
  language=Python,
  basicstyle=\ttfamily\small,
  backgroundcolor=\color{code-bg},
  keywordstyle=\color{suva-blue}\bfseries,
  commentstyle=\color{suva-green}\itshape,
  numbers=left,
  numberstyle=\tiny\color{gray},
  frame=single,
  breaklines=true,
  captionpos=b
}

%--- Code style: bash ---
\lstdefinestyle{bash}{
  language=bash,
  basicstyle=\ttfamily\small,
  backgroundcolor=\color{code-bg},
  keywordstyle=\color{suva-orange}\bfseries,
  commentstyle=\color{gray}\itshape,
  numbers=left,
  numberstyle=\tiny\color{gray},
  frame=single,
  breaklines=true,
  captionpos=b
}

%--- Custom boxes ---
\newtcolorbox{resultbox}[1][]{
  enhanced,colback=suva-green!8,colframe=suva-green!60,
  fonttitle=\bfseries,title=#1,arc=3mm
}
\newtcolorbox{warningbox}[1][]{
  enhanced,colback=suva-orange!8,colframe=suva-orange!60,
  fonttitle=\bfseries,title=#1,arc=3mm
}
\newtcolorbox{infobox}[1][]{
  enhanced,colback=suva-blue!6,colframe=suva-blue!50,
  fonttitle=\bfseries,title=#1,arc=3mm
}

%================================================
\begin{document}
%================================================

%--- Title page ---
\begin{titlepage}
\centering
\vspace*{2cm}
{\Huge\bfseries\textcolor{suva-blue}{SuVa Protocol}\par}
\vspace{0.5cm}
{\Large\textit{Quantum-Resistant Stablecoin on the EVM}\par}
\vspace{1.5cm}
\rule{\linewidth}{0.5pt}
\vspace{0.5cm}
{\LARGE\bfseries Onboarding Guide\par}
\vspace{0.3cm}
{\large From zero to post-quantum: a complete technical introduction\par}
\vspace{0.5cm}
\rule{\linewidth}{0.5pt}
\vspace{2cm}

\begin{tabular}{ll}
\textbf{Project:}   & SuVa Protocol \\[4pt]
\textbf{Author:}    & Mohamed Lamine MBENGUE \\[4pt]
\textbf{Date:}      & April 2026 \\[4pt]
\textbf{Version:}   & 2.0 \\[4pt]
\textbf{Language:}  & English \\[4pt]
\textbf{Codebase:}  & ETHDILITHIUM-main (ZKNox fork) \\
\end{tabular}

\vfill
\begin{tcolorbox}[enhanced,colback=suva-blue!8,colframe=suva-blue,arc=4mm]
\centering\textit{This guide assumes you can write a Python function and have heard\\
of Bitcoin, but know nothing about cryptography or blockchain.\\
By the end, you will understand every line of code in SuVa.}
\end{tcolorbox}
\end{titlepage}

%--- Table of contents ---
\tableofcontents
\newpage

%================================================
\section*{How to use this guide}
\addcontentsline{toc}{section}{How to use this guide}
%================================================

\begin{infobox}[Reading strategy]
This guide is designed to be read from beginning to end, but each part is also self-contained.
If you already know modular arithmetic, start at Part~2.
If you know what a hash function is, jump directly to Section~2.4.
Every mathematical concept is accompanied by a concrete numerical example --- do not skip them.
\end{infobox}

The guide is divided into seven parts:

\begin{enumerate}
  \item \textbf{Mathematics} --- The necessary algebra: modular arithmetic, polynomial rings, Euclidean lattices, and NTT.
  \item \textbf{Cryptography} --- Digital signatures, hash functions, ECDSA, and why quantum computers are dangerous.
  \item \textbf{Post-Quantum Cryptography} --- The new NIST standards: Dilithium and Falcon.
  \item \textbf{Blockchain / EVM} --- Ethereum, gas, Solidity, and Foundry.
  \item \textbf{The SuVa Project} --- Vision, architecture, phases, and key numbers.
  \item \textbf{Code Walkthrough} --- Annotated walkthrough of Python and Solidity code line by line.
  \item \textbf{Practical Guide} --- Set up your environment and run the tests yourself.
\end{enumerate}

\newpage

%================================================
\part{Mathematics}
%================================================

%================================================
\section{Modular Arithmetic}
%================================================

\subsection{What does $a \equiv b \pmod{n}$ mean?}

You already know modular arithmetic without realising it. On a 12-hour clock, if it is 10 o'clock and you add 5 hours, you get 3, not 15. That is arithmetic modulo 12.

Formally, $a \equiv b \pmod{n}$ (read: ``$a$ is congruent to $b$ modulo $n$'') when $a$ and $b$ have the same remainder when divided by $n$, i.e.\ $n$ divides $a - b$.

\begin{infobox}[Definition]
$a \equiv b \pmod{n}$ means $n \mid (a - b)$, i.e.\ there exists an integer $k$ such that $a = b + k \cdot n$.
\end{infobox}

\paragraph{Concrete examples:}
\begin{align*}
17 &\equiv 3 \pmod{7} \quad \text{since } 17 = 2 \times 7 + 3 \\
25 &\equiv 0 \pmod{5} \quad \text{since } 25 = 5 \times 5 + 0 \\
-1 &\equiv 6 \pmod{7} \quad \text{since } -1 = (-1) \times 7 + 6
\end{align*}

The set $\{0, 1, 2, \ldots, n-1\}$ forms what we call $\mathbb{Z}_n$, the integers modulo $n$.

\subsection{Addition and multiplication modulo $n$}

Perform the ordinary operation, then reduce:

\begin{infobox}[Arithmetic in $\mathbb{Z}_7$]
$5 + 4 \equiv 9 \equiv 2 \pmod{7}$ \\
$3 \times 5 \equiv 15 \equiv 1 \pmod{7}$
\end{infobox}

\subsection{Modular inverse}

The \emph{modular inverse} of $a$ modulo $n$ is the integer $a^{-1}$ such that:
\[a \cdot a^{-1} \equiv 1 \pmod{n}\]

\begin{resultbox}[Example: $3^{-1} \pmod{7}$]
We seek $x$ such that $3x \equiv 1 \pmod{7}$.\\
Try $x = 5$: $3 \times 5 = 15 \equiv 1 \pmod{7}$. \checkmark \\
So $3^{-1} \equiv 5 \pmod{7}$.
\end{resultbox}

\begin{warningbox}[When does the inverse exist?]
The inverse of $a$ modulo $n$ exists \emph{if and only if} $\gcd(a, n) = 1$. If $n$ is prime, every non-zero element has an inverse. This is why cryptographic schemes choose $q$ prime.
\end{warningbox}

\subsection{Why modular arithmetic in cryptography?}

\begin{enumerate}
  \item \textbf{Finite and efficient}: all numbers stay in $\{0, \ldots, n-1\}$, so operations are fast even for huge $n$ (e.g.\ $n = 8\,380\,417$).
  \item \textbf{Hard to reverse}: multiplication modulo $n$ is easy in one direction, but recovering the inputs from the result (the \emph{discrete logarithm problem}) is computationally hard.
  \item \textbf{Rich algebraic structure}: $\mathbb{Z}_p$ for prime $p$ forms a \emph{field} --- every non-zero element has an inverse.
\end{enumerate}

%================================================
\section{Polynomial Rings}
%================================================

\subsection{What is a polynomial?}

A polynomial over $\mathbb{Z}_q$ of degree less than $n$ is:
\[p(X) = a_0 + a_1 X + a_2 X^2 + \cdots + a_{n-1} X^{n-1}\]
where each $a_i \in \{0, 1, \ldots, q-1\}$. Think of it as a list of $n$ integers.

\begin{lstlisting}[style=python]
# Polynomial 3 + 2X + 5X^2 + 1X^3  (n=4, coefficients mod 7)
p = [3, 2, 5, 1]   # p[i] = coefficient of X^i
\end{lstlisting}

\subsection{Reduction modulo $(X^n + 1)$: the rule $X^n \equiv -1$}

In the ring $\mathbb{Z}_q[X]/(X^n+1)$, whenever we see $X^n$, we replace it with $-1$:

\begin{infobox}[The reduction rule]
$X^n \equiv -1$, so $X^{n+k} \equiv -X^k$ for all $k \geq 0$.
\end{infobox}

During multiplication, any term $a_k X^k$ with $k \geq n$ becomes $-a_{k-n} X^{k-n}$.

\subsection{Polynomial multiplication in $\mathbb{Z}_q[X]/(X^n+1)$}

\begin{lstlisting}[style=python]
def poly_mul(a, b, q, n):
    result = [0] * n
    for i in range(n):
        for j in range(n):
            idx = (i + j) % n
            prod = (a[i] * b[j]) % q
            if i + j >= n:   # X^(i+j) = -X^(i+j-n)
                result[idx] = (result[idx] - prod) % q
            else:
                result[idx] = (result[idx] + prod) % q
    return result
\end{lstlisting}

\subsection{Why this ring? Efficiency + Security}

\begin{enumerate}
  \item \textbf{Efficiency}: the NTT reduces polynomial multiplication from $O(n^2)$ to $O(n \log n)$.
  \item \textbf{Security}: Ring-LWE problems on this ring are believed to be hard even for quantum computers.
\end{enumerate}

%================================================
\section{Euclidean Lattices}
%================================================

\subsection{What is a lattice?}

A lattice is a regular grid of points in $n$-dimensional space. Given $n$ linearly independent basis vectors $\mathbf{b}_1, \ldots, \mathbf{b}_n$:
\[\mathcal{L} = \left\{ \sum_{i=1}^n z_i \mathbf{b}_i \; : \; z_i \in \mathbb{Z} \right\}\]

\subsection{The Shortest Vector Problem (SVP)}

\begin{warningbox}[Shortest Vector Problem]
Given a basis of a lattice, find the shortest non-zero vector in $\mathcal{L}$.
\end{warningbox}

SVP is NP-hard in the worst case. Even quantum computers only offer a modest speedup (square root via Grover, not polynomial like Shor).

\subsection{Learning With Errors (LWE)}

\begin{infobox}[Learning With Errors]
Given a random matrix $\mathbf{A} \in \mathbb{Z}_q^{m \times n}$ and a vector $\mathbf{b} = \mathbf{A}\mathbf{s} + \mathbf{e}$ where $\mathbf{s}$ is a secret and $\mathbf{e}$ a small error vector, recover $\mathbf{s}$.
\end{infobox}

\paragraph{Analogy:} Noisy GPS coordinates. With exact coordinates you can locate someone; with noise it becomes very hard.

\subsection{The NTRU lattice (foundation of Falcon)}

\begin{infobox}[NTRU Problem]
Given $h \equiv g \cdot f^{-1} \pmod{q, X^n+1}$ with $f$ and $g$ small polynomials, recover $f$ and $g$.
\end{infobox}

\begin{resultbox}[Why lattices are post-quantum]
The best known quantum algorithm for lattice problems runs in time $2^{O(n)}$. No polynomial-time quantum algorithm for SVP or LWE is known or expected.
\end{resultbox}

%================================================
\section{Number Theoretic Transform (NTT)}
%================================================

\subsection{Why NTT?}

Naively multiplying two degree-$n$ polynomials requires $n^2$ multiplications. For $n=256$, that is $65\,536$ multiplications --- expensive in gas. The NTT reduces this to $O(n \log n)$.

\subsection{Principle}

The NTT evaluates a polynomial at $n$ special points (roots of unity modulo $q$), multiplies pointwise, then applies the inverse NTT.

\begin{center}
\begin{tabular}{lll}
\toprule
\textbf{Scheme} & \textbf{Modulus $q$} & \textbf{NTT root} \\
\midrule
Dilithium & $q = 8\,380\,417$ & $\omega = 1753$ \\
Falcon    & $q = 12\,289$     & primitive root mod $q$ \\
\bottomrule
\end{tabular}
\end{center}

\begin{warningbox}[ETHFalcon Phase 2 without NTT]
Phase 2 uses the schoolbook $O(n^2)$ algorithm: 46.67M gas. The NTT would bring this down to $\sim$4M gas.
\end{warningbox}

\newpage
%================================================
\part{Cryptography}
%================================================

%================================================
\section{Digital Signatures --- The Basics}
%================================================

\subsection{What problem do signatures solve?}

Digital signatures guarantee three things:
\begin{itemize}
  \item \textbf{Authenticity}: only the holder of the private key could have signed.
  \item \textbf{Non-repudiation}: the signer cannot deny having signed.
  \item \textbf{Integrity}: if the message is altered after signing, the signature becomes invalid.
\end{itemize}

\subsection{The KeyGen $\to$ Sign $\to$ Verify model}

\begin{itemize}
  \item $\mathbf{KeyGen}()$: generates a \emph{public key} $pk$ (shared) and a \emph{private key} $sk$ (secret).
  \item $\mathbf{Sign}(sk, msg)$: produces a signature $\sigma$.
  \item $\mathbf{Verify}(pk, msg, \sigma)$: returns \texttt{true} or \texttt{false}.
\end{itemize}

%================================================
\section{Hash Functions}
%================================================

\subsection{Security properties}

\begin{enumerate}
  \item \textbf{Preimage resistance}: given $h$, it is infeasible to find $m$ such that $H(m) = h$.
  \item \textbf{Second preimage resistance}: given $m$, it is infeasible to find $m' \neq m$ with $H(m) = H(m')$.
  \item \textbf{Collision resistance}: it is infeasible to find any pair $m \neq m'$ with $H(m) = H(m')$.
\end{enumerate}

\begin{center}
\begin{tabular}{llll}
\toprule
\textbf{Function} & \textbf{Output} & \textbf{Standard} & \textbf{SuVa usage} \\
\midrule
SHA-256    & 256 bits (fixed) & NIST FIPS 180 & Not directly \\
Keccak-256 & 256 bits (fixed) & Ethereum native & ETHFalcon HashToPoint \\
SHAKE-256  & Variable (XOF)   & NIST FIPS 202  & Standard Falcon/Dilithium \\
\bottomrule
\end{tabular}
\end{center}

\subsection{Why Keccak is native on the EVM}

The EVM has a native \texttt{KECCAK256} opcode: 30 gas + 6 gas per 32-byte word. SHAKE-256 has no native EVM opcode --- implementing it in Solidity would cost millions of gas.

\begin{resultbox}[SuVa's innovation]
ETHFalcon replaces SHAKE-256 with a deterministic PRNG built on Keccak (KeccakPRNG). This single change reduces the on-chain verification cost by an order of magnitude while preserving the lattice security.
\end{resultbox}

%================================================
\section{ECDSA --- Ethereum's Current Signature}
%================================================

\subsection{Private key $\to$ Public key}

\begin{itemize}
  \item \textbf{Private key}: a random integer $k \in \{1, \ldots, n-1\}$ (256 bits)
  \item \textbf{Public key}: $K = k \cdot G$ where $G$ is the secp256k1 generator point
\end{itemize}

Computing $k \cdot G$ is easy (scalar multiplication). Recovering $k$ from $K$ is the ECDLP --- believed computationally infeasible classically.

%================================================
\section{Why Quantum Computers Break ECDSA}
%================================================

\begin{warningbox}[Shor's Algorithm (1994)]
A quantum computer can solve the discrete logarithm problem (ECDLP) in \textbf{polynomial time} $O(n^3)$. With $\sim$4000 stable qubits, ECDSA is completely broken.
\end{warningbox}

\begin{warningbox}[``Harvest Now, Decrypt Later'' attack]
An adversary can \textbf{today} record all public Ethereum transactions on the blockchain. When a quantum computer arrives, they will recover private keys from historical public keys. Migration to PQC is urgent \emph{now}, not only when quantum computers arrive.
\end{warningbox}

\begin{infobox}[Grover vs Shor]
\textbf{Shor}: polynomial time --- completely breaks RSA/ECDSA.\\
\textbf{Grover}: square-root speedup --- halves the security of hash functions. SHA-256 remains safe (128 quantum bits). Hash-based signatures (SPHINCS+) remain secure.
\end{infobox}

\newpage
%================================================
\part{Post-Quantum Cryptography}
%================================================

%================================================
\section{NIST PQC Standardisation}
%================================================

In 2016, NIST launched a public competition. After 8 years and 3 evaluation rounds, four standards were published in August 2024:

\begin{center}
\begin{tabular}{llll}
\toprule
\textbf{FIPS} & \textbf{Algorithm} & \textbf{Type} & \textbf{Hard problem} \\
\midrule
FIPS 203 & ML-KEM (Kyber)     & Key encapsulation  & Module-LWE \\
FIPS 204 & ML-DSA (Dilithium) & Digital signature  & Module-LWE + SIS \\
FIPS 205 & SLH-DSA (SPHINCS+) & Digital signature  & Hash functions \\
FIPS 206 & FN-DSA (Falcon)    & Digital signature  & NTRU lattice \\
\bottomrule
\end{tabular}
\end{center}

SuVa uses ML-DSA (Phase 1), FN-DSA/Falcon (Phases 2--4), and will study SLH-DSA/SPHINCS+ in Phase 5.

%================================================
\section{Falcon (FN-DSA, FIPS 206)}
%================================================

\subsection{Parameters}

\begin{center}
\begin{tabular}{lll}
\toprule
\textbf{Parameter} & \textbf{Falcon-256} & \textbf{Falcon-512} \\
\midrule
$n$ & 256 & 512 \\
$q$ & 12\,289 & 12\,289 \\
Public key size & 448 bytes & 897 bytes \\
Signature size  & 356 bytes & 666 bytes \\
\bottomrule
\end{tabular}
\end{center}

SuVa uses Falcon-256 ($n=256$). Signatures are 3.6$\times$ shorter than Dilithium.

\subsection{Verification: step by step}

\begin{infobox}[ETHFalcon verification algorithm]
\textbf{Inputs}: public key $h$, salt (40 bytes), $s_1$ (256 polynomial coefficients), message $msg$

\begin{enumerate}
  \item \textbf{HashToPoint}: $t \gets \text{KeccakPRNG}(\text{salt} \| msg)$ --- produces 256 uniform coefficients in $\{0,\ldots,q-1\}$
  \item \textbf{Polynomial multiplication}: $h \cdot s_1 \pmod{q, X^{256}+1}$
  \item \textbf{Recover} $s_0$: $s_0 \equiv t - h \cdot s_1 \pmod{q}$
  \item \textbf{Centre} $s_0$: map each coefficient from $\{0,\ldots,q-1\}$ to $(-q/2, q/2]$
  \item \textbf{Norm check}: $\|s_0\|^2 + \|s_1\|^2 \leq \beta^2$
\end{enumerate}

Key idea: a valid signature has small $(s_0, s_1)$. Forging without the trapdoor is computationally infeasible.
\end{infobox}

%================================================
\section{ETHFalcon --- SuVa's Innovation}
%================================================

\subsection{KeccakPRNG --- HashToPoint}

\begin{lstlisting}[style=python, caption={KeccakPRNG for HashToPoint}]
def hash_to_point_keccak(salt, msg, n, q):
    state = keccak256(salt + msg)   # initial seed
    result = []
    counter = 0
    while len(result) < n:
        block = keccak256(state + counter.to_bytes(8, 'big'))
        for i in range(0, 32, 2):
            elt = (block[i] << 8) | block[i+1]
            if elt < 5 * q:       # rejection for uniform distribution
                result.append(elt % q)
        counter += 1
    return result[:n]
\end{lstlisting}

The threshold $5 \times 12\,289 = 61\,445$ guarantees a perfectly uniform distribution in $\{0,\ldots,q-1\}$.

\begin{resultbox}[Security argument]
The security of Falcon's verification equation does not depend on the specific hash function used to generate $t$, only on: (1) $t$ is uniformly distributed; (2) $t$ is determined by $(salt, msg)$; (3) the function is preimage- and collision-resistant. Keccak-256 satisfies all three. The lattice hardness argument is unchanged.
\end{resultbox}

\newpage
%================================================
\part{Blockchain and EVM}
%================================================

%================================================
\section{Ethereum Basics}
%================================================

\subsection{Accounts: EOA vs Contract Accounts}

\begin{center}
\begin{tabular}{lll}
\toprule
\textbf{Property} & \textbf{EOA} & \textbf{Contract Account} \\
\midrule
Controlled by & ECDSA private key & Smart contract code \\
Has code?     & No  & Yes \\
Can initiate tx? & Yes & No (only via calls) \\
\bottomrule
\end{tabular}
\end{center}

SuVa Phase 4 makes contract accounts act as wallets where the ``private key'' is a Falcon key instead of ECDSA.

\subsection{Gas}

Every EVM operation has a cost measured in \textbf{gas}. You pay: $\text{fee} = \text{gas used} \times \text{gas price (gwei)}$.

\begin{center}
\begin{tabular}{lr}
\toprule
\textbf{Operation} & \textbf{Gas cost} \\
\midrule
ADD, SUB                & 3 \\
MUL                     & 5 \\
MULMOD                  & 8 \\
KECCAK256 (32 bytes)    & 36 \\
SLOAD (cold storage)    & 2\,100 \\
Base transaction        & 21\,000 \\
Block gas limit (mainnet) & 30\,000\,000 \\
\bottomrule
\end{tabular}
\end{center}

%================================================
\section{Key Solidity Concepts}
%================================================

\subsection{Contract structure}

\begin{lstlisting}[style=solidity, caption={Key Solidity elements}]
// SPDX-License-Identifier: MIT
pragma solidity ^0.8.25;

contract MyContract {
    // State variables (on-chain, expensive)
    uint256 public counter;
    address public owner;

    // Modifier: reusable condition check
    modifier onlyOwner() {
        require(msg.sender == owner, "Access denied");
        _;
    }

    // pure = no state read/write
    function add(uint256 a, uint256 b)
        public pure returns (uint256)
    {
        return a + b;
    }

    // view = reads state, no write
    function getCounter() public view returns (uint256) {
        return counter;
    }

    // payable = can receive ETH
    function deposit() public payable { }
}
\end{lstlisting}

\subsection{Data locations}

\begin{center}
\begin{tabular}{llr}
\toprule
\textbf{Location} & \textbf{Description} & \textbf{Relative cost} \\
\midrule
\texttt{calldata} & Read-only function input & Cheapest \\
\texttt{memory}   & Temporary, cleared after call & Medium \\
\texttt{storage}  & Persistent on-chain state & Most expensive \\
\bottomrule
\end{tabular}
\end{center}

In ETHFalcon, $h$, $salt$, $s_1$, $msg$ are declared \texttt{calldata} to avoid unnecessary copies.

\subsection{Foundry: compile, test, deploy}

\begin{lstlisting}[style=bash]
# Compile contracts
forge build

# Run all tests with detail
forge test -vv

# Tests with gas report
forge test -vv --gas-report

# Deploy to Sepolia
forge script script/DeploySuVa.s.sol \
  --rpc-url https://ethereum-sepolia-rpc.publicnode.com \
  --private-key YOUR_PRIVATE_KEY \
  --broadcast

# Verify a contract on Etherscan
forge verify-contract ADDRESS src/Contract.sol:Contract \
  --chain-id 11155111 \
  --etherscan-api-key YOUR_API_KEY \
  --watch
\end{lstlisting}

\newpage
%================================================
\part{The SuVa Project}
%================================================

%================================================
\section{Vision and Global Architecture}
%================================================

\subsection{Problem solved}

Financial institutions want to store stablecoins (USDT, USDC) securely on Ethereum. But ECDSA exposes them to the future quantum threat. SuVa adds a post-quantum signature layer without modifying Ethereum.

\subsection{The four phases}

\begin{longtable}{@{}p{1.8cm}p{4cm}p{4cm}p{4cm}@{}}
\toprule
\textbf{Phase} & \textbf{Objective} & \textbf{Deliverable} & \textbf{Tests} \\
\midrule
\endhead
Phase 1 & Dilithium verifier on-chain & \texttt{ZKNOX\_dilithium.sol} & -- \\[4pt]
Phase 2 & ETHFalcon (Falcon + Keccak) & \texttt{ZKNOX\_ethfalcon.sol} + Python & -- \\[4pt]
Phase 3 & QuantumVault & \texttt{QuantumVault.sol}, \texttt{qUSDT.sol} & 26/26 \\[4pt]
Phase 4 & SuVaWallet ERC-4337 & \texttt{SuVaWallet.sol}, \texttt{SuVaWalletFactory.sol} & 15/15 \\[4pt]
Phase 5 & SPHINCS+ PoC & Feasibility study + L2 PoC & upcoming \\
\bottomrule
\end{longtable}

\textbf{Total: 41/41 tests passing. All contracts deployed and verified on Sepolia.}

%================================================
\section{Phase 3 --- QuantumVault}
%================================================

\subsection{What is QuantumVault?}

QuantumVault is a stablecoin vault secured by Falcon. A user deposits USDT, receives qUSDT (tokenised receipt), and can only withdraw by presenting a valid Falcon signature.

\subsection{Contract architecture}

\begin{center}
\begin{tabular}{ll}
\toprule
\textbf{Contract} & \textbf{Role} \\
\midrule
\texttt{ZKNOX\_ethfalcon.sol} & On-chain Falcon verifier (Phase 2) \\
\texttt{MockERC20.sol}        & Simulated USDT for Sepolia \\
\texttt{qUSDT.sol}            & ERC-20 receipt token for deposits \\
\texttt{QuantumVault.sol}     & Main vault contract \\
\bottomrule
\end{tabular}
\end{center}

\subsection{Simplified interface}

\begin{lstlisting}[style=solidity, caption={QuantumVault interface}]
// Register your Falcon public key
function registerFalconKey(uint256[256] calldata h) external;

// Deposit USDT -> receive qUSDT
function deposit(uint256 amount) external;

// Withdraw: must present a valid Falcon signature
function withdraw(
    uint256 amount,
    uint256[256] calldata h,
    bytes   calldata salt,
    uint256[256] calldata s1
) external;
\end{lstlisting}

\subsection{Phase 3 key figures}

\begin{center}
\begin{tabular}{lr}
\toprule
\textbf{Metric} & \textbf{Value} \\
\midrule
Forge tests & 26/26 passing \\
Gas deposit & $\sim$50\,000 gas \\
Gas withdraw (with Falcon verify) & $\sim$48M gas (schoolbook) \\
Sepolia address QuantumVault & \texttt{0x601c...af1c4} \\
\bottomrule
\end{tabular}
\end{center}

%================================================
\section{Phase 4 --- SuVaWallet ERC-4337}
%================================================

\subsection{What is ERC-4337?}

\textbf{The problem:} A standard Ethereum account (EOA) is tied to an ECDSA key. Using Falcon natively is impossible.

\textbf{The solution (ERC-4337):} A smart contract can act as an account if it implements \texttt{validateUserOp()}.

\begin{center}
\begin{tabular}{ll}
\toprule
\textbf{Component} & \textbf{Role} \\
\midrule
\texttt{UserOperation} & The user's ``transaction'' \\
\texttt{EntryPoint} & Official ERC-4337 contract (\texttt{0x5FF1...2789}) \\
\texttt{Bundler} & Batches UserOps and calls EntryPoint \\
\texttt{SuVaWallet} & Our contract-account \\
\bottomrule
\end{tabular}
\end{center}

\subsection{Transaction flow}

\begin{lstlisting}
User --> signs with ECDSA + Falcon
  --> sends UserOperation to Bundler
  --> Bundler calls EntryPoint
  --> EntryPoint calls SuVaWallet.validateUserOp()
      --> verifies ECDSA (ecrecover)
      --> verifies Falcon (IFalconVerifier)
      --> returns SIG_VALIDATION_SUCCESS (0)
  --> EntryPoint calls SuVaWallet.execute()
      --> executes the transaction
\end{lstlisting}

\subsection{Dual signature format}

\begin{lstlisting}[style=solidity, caption={UserOperation signature layout}]
// bytes[0..64]  = ECDSA signature (65 bytes)
// bytes[65..]   = abi.encode(h, salt, s1) -- Falcon

bytes memory ecdsaSig  = sig[0:65];
bytes memory falconSig = sig[65:];

(uint256[256] memory h,
 bytes        memory salt,
 uint256[256] memory s1) = abi.decode(falconSig, (...));
\end{lstlisting}

Both signatures must be valid for the transaction to execute.

\subsection{SuVaWalletFactory --- CREATE2}

\begin{lstlisting}[style=solidity]
function createWallet(
    address entryPoint,
    address falconVerifier,
    address owner,
    uint256 salt
) external returns (SuVaWallet wallet) {
    bytes32 create2salt = keccak256(abi.encodePacked(owner, salt));
    wallet = new SuVaWallet{salt: create2salt}(
        entryPoint, falconVerifier, owner
    );
}
\end{lstlisting}

\subsection{Public interface --- IFalconVerifier}

\begin{lstlisting}[style=solidity, caption={Shared interface for interoperability}]
interface IFalconVerifier {
    function verify(
        uint256[256] calldata h,       // public key
        bytes        calldata salt,    // signature nonce
        uint256[256] calldata s1,      // vector s1
        bytes        calldata message  // signed message
    ) external pure returns (bool);
}
\end{lstlisting}

This interface is the interoperability point with BTQ (see Part~7).

\subsection{Sepolia addresses --- Phase 4}

\begin{center}
\begin{tabular}{ll}
\toprule
\textbf{Contract} & \textbf{Sepolia Address} \\
\midrule
SuVaWalletFactory & \texttt{0x716a0C550eC1267Db493515d4a114C2FCf942e6B} \\
SuVaWallet & \texttt{0x8a907C21014ED2926a53b4d18eC6e1b8C5a514cE} \\
EntryPoint (ERC-4337) & \texttt{0x5FF137D4b0FDCD49DcA30c7CF57E578a026d2789} \\
\bottomrule
\end{tabular}
\end{center}

Both contracts are \textbf{verified on Sepolia Etherscan} (source code publicly readable).

%================================================
\section{EIP-8141 and Phase 5}
%================================================

\subsection{EIP-8141 --- Frame Transaction (April 2026)}

EIP-8141, published by Vitalik Buterin and the ERC-4337 team, proposes integrating account abstraction \textbf{natively into the Ethereum protocol}:

\begin{center}
\begin{tabular}{lll}
\toprule
& \textbf{ERC-4337 (today)} & \textbf{EIP-8141 (future)} \\
\midrule
Level & Application & Native protocol \\
Bundler & Required & Not required \\
EntryPoint & External contract & Built into EVM \\
PQC support & Via \texttt{validateUserOp()} & Native \\
\bottomrule
\end{tabular}
\end{center}

\textbf{SuVa Phase 4 is a direct precursor of this vision.}

\subsection{Phase 5 --- SPHINCS+}

SPHINCS+ (FIPS 205, SLH-DSA) is a hash-based signature scheme --- no algebraic structure, no lattices, no curves. Security relies only on the hardness of SHA-256/SHAKE, the most conservative assumption possible.

\textbf{Phase 5 goal:} Feasibility study and PoC on L2, in the direction of EIP-8141.

\newpage
%================================================
\part{Code Walkthrough}
%================================================

%================================================
\section{Source Code Structure}
%================================================

\begin{lstlisting}
d:\Projects\suva\
|- ONBOARDING.md
|- rapports\
|   |- rapport_phase3_EN.tex  / rapport_phase3_FR.tex
|   |- rapport_phase4_EN.tex  / rapport_phase4_FR.tex
|   |- ONBOARDING_EN.tex      / ONBOARDING_FR.tex
`- falcon_dilithium\
    `- ETHDILITHIUM-main\
        |- src\
        |   |- IFalconVerifier.sol    <- shared interface
        |   |- ZKNOX_ethfalcon.sol    <- Falcon verifier
        |   |- QuantumVault.sol       <- Phase 3
        |   |- qUSDT.sol              <- Phase 3
        |   |- MockERC20.sol          <- Phase 3
        |   |- SuVaWallet.sol         <- Phase 4
        |   `- SuVaWalletFactory.sol  <- Phase 4
        |- test\
        |   |- QuantumVault.t.sol     <- 26 tests
        |   `- SuVaWallet.t.sol       <- 15 tests
        |- script\
        |   `- DeploySuVa.s.sol       <- deployment script
        `- python-ref\
            `- Falcon_dilithium_py\
                |- Falcon_falcon\     <- Python ETHFalcon module
                `- generate_falcon_test_vectors.py
\end{lstlisting}

%================================================
\section{Reading SuVaWallet.sol}
%================================================

\begin{lstlisting}[style=solidity, caption={SuVaWallet.sol --- annotated structure}]
contract SuVaWallet is IAccount {
    address public immutable entryPoint;
    address public immutable falconVerifier;
    address public immutable owner;

    // Hash of the registered Falcon public key
    bytes32 public falconPkHash;

    // Only the ERC-4337 EntryPoint can call these functions
    modifier onlyEntryPoint() {
        require(msg.sender == entryPoint, "Not EntryPoint");
        _;
    }

    // UserOperation validation:
    // 1. Verify ECDSA
    // 2. Verify Falcon
    // 3. Return 0 (success) or 1 (failure)
    function validateUserOp(
        UserOperation calldata userOp,
        bytes32 userOpHash,
        uint256 missingAccountFunds
    ) external onlyEntryPoint returns (uint256) {
        _verifyECDSA(userOp.signature, userOpHash);
        _verifyFalcon(userOp.signature, userOpHash);
        return 0;  // SIG_VALIDATION_SUCCESS
    }

    // Execute a transaction after validation
    function execute(
        address to,
        uint256 value,
        bytes calldata data
    ) external onlyEntryPoint {
        (bool success,) = to.call{value: value}(data);
        require(success, "Execute: failed");
    }
}
\end{lstlisting}

%================================================
\section{Forge Tests --- Understanding Mocks}
%================================================

\begin{lstlisting}[style=solidity, caption={SuVaWallet.t.sol --- test structure}]
contract SuVaWalletTest is Test {

    // Mocks simulate real contracts in a controlled environment
    MockEntryPoint     entryPoint;
    MockFalconVerifier falconVerifier;
    SuVaWallet         wallet;

    uint256 constant OWNER_PRIVATE_KEY = 0xA11CE;
    address OWNER;

    function setUp() public {
        OWNER = vm.addr(OWNER_PRIVATE_KEY);
        entryPoint     = new MockEntryPoint();
        falconVerifier = new MockFalconVerifier();
        wallet = new SuVaWallet(
            address(entryPoint),
            address(falconVerifier),
            OWNER
        );
    }

    // Test: validateUserOp must return 0 when both sigs are valid
    function test_ValidateUserOp_Success() public {
        bytes32 hash = keccak256("test");
        (uint8 v, bytes32 r, bytes32 s) =
            vm.sign(OWNER_PRIVATE_KEY, hash);
        bytes memory ecdsaSig = abi.encodePacked(r, s, v);

        // MockFalconVerifier always returns true
        bytes memory falconPart = abi.encode(h, salt, s1);
        bytes memory fullSig = abi.encodePacked(ecdsaSig, falconPart);

        uint256 result = wallet.validateUserOp(userOp, hash, 0);
        assertEq(result, 0, "Must return SIG_VALIDATION_SUCCESS");
    }
}
\end{lstlisting}

\newpage
%================================================
\part{Practical Guide}
%================================================

%================================================
\section{Environment Setup}
%================================================

\begin{lstlisting}[style=bash]
# 1. Install Foundry (forge, cast, anvil)
curl -L https://foundry.paradigm.xyz | bash
foundryup

# 2. Verify installation
forge --version

# 3. Install Python dependencies
cd falcon_dilithium/ETHDILITHIUM-main/python-ref/
pip install falcon-py pytest
\end{lstlisting}

%================================================
\section{Running the Tests}
%================================================

\begin{lstlisting}[style=bash]
cd falcon_dilithium/ETHDILITHIUM-main

# Compile
forge build

# Fast tests (lite profile, no optimiser)
FOUNDRY_PROFILE=lite forge test -vv

# Full tests with gas report
forge test -vv --gas-report

# Expected output:
# [PASS] test_ValidateUserOp_Success
# [PASS] test_ValidateUserOp_WrongECDSA
# ...
# Suite result: ok. 41 passed; 0 failed
\end{lstlisting}

%================================================
\section{Deploy and Verify}
%================================================

\begin{lstlisting}[style=bash]
# Deploy to Sepolia
forge script script/DeploySuVa.s.sol \
  --rpc-url https://ethereum-sepolia-rpc.publicnode.com \
  --private-key $PRIVATE_KEY \
  --broadcast

# Verify SuVaWalletFactory
forge verify-contract \
  0x716a0C550eC1267Db493515d4a114C2FCf942e6B \
  src/SuVaWalletFactory.sol:SuVaWalletFactory \
  --chain-id 11155111 \
  --etherscan-api-key $ETHERSCAN_KEY \
  --watch

# Verify SuVaWallet (with constructor args)
forge verify-contract \
  0x8a907C21014ED2926a53b4d18eC6e1b8C5a514cE \
  src/SuVaWallet.sol:SuVaWallet \
  --chain-id 11155111 \
  --constructor-args $(cast abi-encode \
    "constructor(address,address,address)" \
    0x5FF137D4b0FDCD49DcA30c7CF57E578a026d2789 \
    0xa6279c2ce813306da0f7b81e434e4a88f0b25626 \
    0xD7390D91139Bf5868DDdF3b6677708b30b78D210) \
  --etherscan-api-key $ETHERSCAN_KEY \
  --watch
\end{lstlisting}

%================================================
\section{Reference Addresses --- Sepolia}
%================================================

\begin{longtable}{@{}p{2cm}p{4cm}p{8cm}@{}}
\toprule
\textbf{Phase} & \textbf{Contract} & \textbf{Sepolia Address} \\
\midrule
\endhead
2 / 3 & ZKNOX\_ethfalcon & \texttt{0xa6279c2ce813306da0f7b81e434e4a88f0b25626} \\
3 & MockUSDT & \texttt{0xf088bc3f822ee138753afe3190d1adec60abbaf3} \\
3 & qUSDT & \texttt{0xf8f3009c473d538da9b188f33bd6a6cbead29902} \\
3 & QuantumVault & \texttt{0x601cd837e9edd0dfb311ff1a48ed7bea7bbaf1c4} \\
4 & SuVaWalletFactory & \texttt{0x716a0C550eC1267Db493515d4a114C2FCf942e6B} \\
4 & SuVaWallet & \texttt{0x8a907C21014ED2926a53b4d18eC6e1b8C5a514cE} \\
-- & EntryPoint (ERC-4337) & \texttt{0x5FF137D4b0FDCD49DcA30c7CF57E578a026d2789} \\
\bottomrule
\end{longtable}

%================================================
\section{BTQ Collaboration and Future Direction}
%================================================

\subsection{BTQ and QSSN}

BTQ is a post-quantum blockchain startup. Their QSSN (Quantum Safe Signature Network) protocol shares the same goal as SuVa: quantum-safe assets on Ethereum without modifying the protocol.

\textbf{Points of convergence:}
\begin{itemize}
  \item Same approach: wrapping existing assets with PQC signatures
  \item Shared interface: \texttt{IFalconVerifier} can become a standard
  \item Joint EIP: submit a quantum-safe Ethereum standard for the whole community
\end{itemize}

\subsection{Glossary}

\begin{longtable}{@{}p{3cm}p{11cm}@{}}
\toprule
\textbf{Term} & \textbf{Definition} \\
\midrule
\endhead
EOA & Standard Ethereum account linked to an ECDSA key \\
Smart contract & Program running on the Ethereum blockchain \\
Gas & Computation cost of an Ethereum operation (paid in ETH) \\
ERC-20 & Standard for fungible tokens (USDT, USDC, etc.) \\
ERC-4337 & Standard for contract-accounts with custom validation \\
Lattice & Mathematical structure used by Falcon (Euclidean networks) \\
NTT & Number Theoretic Transform --- accelerates polynomial multiplication \\
Bundler & Batches UserOps and submits them to the EntryPoint \\
Sepolia & Ethereum test network (free, no real money) \\
Etherscan & Blockchain explorer (source code, transactions, addresses) \\
CREATE2 & Deterministic contract deployment (same address every time) \\
SPHINCS+ & Hash-based PQC signature algorithm (Phase 5) \\
QSSN & Quantum Safe Signature Network --- BTQ's protocol \\
EIP & Ethereum Improvement Proposal \\
EIP-8141 & Native account abstraction proposal for Ethereum (2026) \\
SVP & Shortest Vector Problem --- hard lattice problem \\
LWE & Learning With Errors --- Dilithium's security foundation \\
NTRU & Lattice problem --- Falcon's security foundation \\
XOF & Extendable Output Function (SHAKE-256) \\
\bottomrule
\end{longtable}

\vspace{2em}
\begin{center}
\textit{SuVa Protocol --- Quantum-Safe Ethereum, Today.}
\end{center}

\end{document}
